% ---------------------------------------------------------------------
%  01 · Contexto
%  Ideas de la presentación original (nodo heterogéneo, potencia y energía,
%  costo y disponibilidad, DVFS depende de la fase). Un ícono grande y una
%  frase corta por columna; el resto, de viva voz.
% ---------------------------------------------------------------------
\seccion{Contexto}

% \ctxcol{ícono}{título}{frase}  (columna de ancho fijo, tarjeta con ícono grande)
\newcommand{\ctxcol}[3]{%
  \begin{minipage}[t]{0.3\textwidth}
    \centering
    \tikz[baseline=0]\node[circle,fill=uisgreen,text=white,minimum size=17mm,inner sep=0pt,
      font=\fontsize{20}{22}\selectfont]{#1};\par\vspace{3mm}
    {\large\bfseries\color{uisgreen}#2\par}\vspace{1.5mm}
    {\small #3\par}
  \end{minipage}}

\begin{frame}{La energía limita cuánto puede escalar un nodo CPU--GPU}
  \vspace{4mm}
  \centering
  \ctxcol{\faMicrochip}{Nodo heterogéneo}{CPU y GPU comparten la carga científica}\hfill
  \ctxcol{\faBolt}{Potencia y energía}{Cada dispositivo consume según su frecuencia y su actividad}\hfill
  \ctxcol{\faLeaf}{Costo y disponibilidad}{La energía incide en los costos operativos y en la sostenibilidad}

  \vspace{7mm}
  \begin{cardhi}[halign=center,fontupper=\normalsize,top=3mm,bottom=3mm,width=0.88\textwidth,center]
    DVFS ajusta voltaje y frecuencia en el hardware, pero\\
    \textbf{el beneficio depende de lo que hace la aplicación en cada fase}
  \end{cardhi}
\end{frame}
